% Generated draft: requires mathematical and source review.
\documentclass[11pt]{article}
\usepackage[margin=25mm]{geometry}
\usepackage{fontspec}
\setmainfont{Latin Modern Roman}
\usepackage{amsmath,amssymb,mathtools,unicode-math}
\setmathfont{Latin Modern Math}
\usepackage{enumitem,needspace,xurl,hyperref}
\setlength{\parindent}{0pt}
\setlength{\parskip}{5pt}
\setlength{\emergencystretch}{3em}
\newcommand{\N}{\mathbb{N}}
\newcommand{\Z}{\mathbb{Z}}
\newcommand{\Q}{\mathbb{Q}}
\newcommand{\R}{\mathbb{R}}
\newcommand{\C}{\mathbb{C}}
\newcommand{\F}{\mathbb{F}}
\newcommand{\A}{\mathbb{A}}
\newcommand{\PP}{\mathbb{P}}
\newcommand{\E}{\mathbb{E}}
\newcommand{\eps}{\varepsilon}
\newcommand{\dd}{\,\mathrm{d}}
\begin{document}
\section*{UnsolvedMath ID: 30000999}
\textbf{Source problem code:} catalogue entry
\subsection{Definitions and mathematical statement}
Let $S^{n-1}\subset\R^n$ be the unit sphere. For a fixed finite $p\geq1$ and a specified metric $d$ on the sphere, the Wasserstein distance between probability measures is
\[
d_{W_p}(\mu,\nu)=
\left(\inf_{\pi\in\Pi(\mu,\nu)}\int d(x,y)^p\,d\pi(x,y)\right)^{1/p},
\]
where $\Pi(\mu,\nu)$ is the set of couplings. The geodesic metric is a conventional choice, but the record does not specify the metric, the range of $p$ or the dimension range.

Let $\mathcal R$ be the source's geodesic Radon transform, described as sending a probability measure to its equator averages. A conventional probability-valued equator-averaging interpretation, for $n\geq2$, uses the normalized measure $\sigma_x$ on $E_x=\{y\in S^{n-1}:x\cdot y=0\}$ and defines
\[
(\mathcal R\mu)(A)=\int_{S^{n-1}}\sigma_x(A\cap E_x)\,d\mu(x).
\]
This interpretation is conditional; the record gives no definition on arbitrary singular measures or precise output space. If the output is instead a different space of slices or measures, its metric must be specified before using Wasserstein distance there.

The literal source question is whether, for each intended $(p,n)$, there is a finite constant $C'(p,n)$ such that
\[
d_{W_p}(\mu,\nu)\leq C'(p,n)\,
d_{W_p}(\mathcal R\mu,\mathcal R\nu)
\quad\text{for all probability measures }\mu,\nu\text{ on }S^{n-1}.
\]
The constant is independent of both measures. No smoothness, density, positivity-bound or antipodal-symmetry restriction on them is stated. A reverse estimate of this form requires injectivity on the stated domain; under ordinary equator averaging, $E_x=E_{-x}$ creates an antipodal identification. The source must therefore clarify the intended transform and any symmetry or quotient convention. The unrestricted quantifier is retained, not silently replaced by one over even measures or a projective quotient.

The corrected-verified note says the mathematical content was checked and normalized. The dated assessment checked 2026-08-21 instead describes only the fragment ``how does this affect the slice averages?'' and says no complete proposition was extracted. That explanation conflicts with the full inequality now supplied. The registry's open label is retained as source metadata, not treated as a verified current status of a fully specified transform.
\subsection{Short English statement}
Ask for a measure-independent inverse Wasserstein bound for spherical equator averaging. Preserve the source's all-probability-measures quantifier while flagging the missing transform definition and antipodal injectivity issue, as well as its stale fragment-only assessment.
\subsection{Sources}
\begin{itemize}
\item Local UnsolvedMath registry, permanent ID 30000999, source code catalogue entry; supplied complete source record, including the corrected-verified statement note and conflicting fragment-only dated assessment.
\item Calculus of Variations (2009), Oberwolfach Reports, DOI 10.4171/owr/2008/31, pp. 1707--1770: \url{https://doi.org/10.4171/owr/2008/31}. Proposed-year metadata gives 2008.
\end{itemize}
\end{document}
